
\subsection{Deep Anomaly Detection in Multivariate Time Series}

Multivariate time-series anomaly detection has evolved from classical statistical methods toward deep architectures that learn normal behavioural patterns and identify deviations~\cite{pang2021deep_survey}. Foundational sequence models LSTM~\cite{hochreiter1997lstm}, GRU~\cite{cho2014gru}, and the TCN~\cite{bai2018tcn} established the architectural vocabulary for temporal modelling. The TCN, in particular, demonstrated that dilated causal convolutions provide a competitive alternative to recurrence, with advantages in training stability and parallel computation. WaveNet~\cite{oord2016wavenet} originally introduced dilated causal convolutions for sequence modelling; subsequent comprehensive reviews~\cite{zhang2020tcn_review} catalogue the diverse extensions that have since been applied to industrial time series.

Recent work has extended TCNs with complementary techniques: TCN-AE architectures~\cite{thill_temporal_2021} leverage autoencoding for unsupervised detection; TCN-Linear hybrids~\cite{wang_tcn-linear_2024} combine convolutional and linear residual paths; and attention-augmented TCNs~\cite{mehta_nac-tcn_2023} incorporate self-attention for industrial monitoring applications. Reconstruction-based anomaly detectors such as USAD~\cite{audibert2020usad} and LSTM-ED~\cite{malhotra2016lstmed} learn normal-behaviour priors and threshold residuals at inference. Transformer-based detectors such as TranAD~\cite{tuli2022tranad} and Anomaly Transformer~\cite{xu2022anomalytransformer} introduce self-attention for anomaly scoring; recent designs (e.g., Informer~\cite{zhou2021informer}, TimesNet~\cite{wu2022timesnet}, MTAD-GAT~\cite{zhao2020mtadgat}) further enrich the attention vocabulary but do not incorporate predictive uncertainty into the attention mechanism itself.

Despite this progress, existing TCN-based detectors treat sensor readings as generic multivariate inputs without modelling sensor-specific reliability characteristics. Our work explicitly couples uncertainty estimation to attention computation and anomaly scoring, adapting model behaviour to the varying reliability of individual sensor channels. Uncertainty quantification (UQ) transforms anomaly detection from threshold-based heuristics into confidence-aware decisions~\cite{kendall2017uncertainties}. Dropout as a Bayesian approximation~\cite{gal2016dropout_bayesian} provides a lightweight mechanism for estimating predictive uncertainty from deterministic architectures, building on the original dropout regularizer~\cite{srivastava2014dropout}; deep ensembles~\cite{lakshminarayanan2017ensembles} offer an alternative with stronger calibration at higher computational cost. MC dropout has been validated across diverse time-series domains: electricity demand forecasting with TCNs~\cite{ghimire_electricity_2025}, stochastic differential equation-based hydrological forecasting~\cite{cui_sde-hnn_2025}, and multivariate imputation. These works demonstrate that MC dropout provides meaningful uncertainty estimates, but they apply uncertainty post hoc as a separate scoring term rather than feeding it back into the model's internal representations. Our framework advances this field of study by closing the loop: uncertainty estimates modulate attention weights and condition the dynamic scoring function, creating integration between UQ and model computation motivated by sensor-specific reliability concerns.

For attention mechanisms in temporal data, the Transformer's self-attention~\cite{vaswani2017attention} enables the direct calculation of relationships between tokens and captures long-range dependencies that might be missed by convolutional filters. Standard attention assumes all input positions are equally trustworthy an assumption violated by physical sensor data where individual measurements carry sensor-specific uncertainty. Conformal prediction~\cite{angelopoulos2021conformal} and concentration-inequality machinery~\cite{boucheron2013concentration} provide distribution-free statistical guarantees, but operate at the output level rather than altering internal representations. Our AUAA mechanism integrates uncertainty directly into attention computation, enabling the model to learn which temporal patterns deserve attention based on the estimated reliability of the underlying sensor measurements; the resulting hybrid score is in turn bounded by a Chebyshev-type inequality (\ref{thm:chebyshev_fp}), providing the framework an explicit distribution-free false-positive guarantee. Multi-task uncertainty weighting~\cite{kendall2018multitask} offers a related precedent for learning task-specific scalings, which our Dynamic Weight Adapter extends into a context-dependent (per-window) scheme. Interpretability remains a critical gap in practical anomaly detection. Methods such as SHAP and LIME provide post-hoc explanations but add computational overhead. Our feature-importance learner is trained jointly with the detection model, producing sensor-specific attributions as a native output without additional post-processing.

\subsection{Benchmarks and Cross-Dataset Evaluation for Time-Series Anomaly Detection}

The benchmarking practice in multivariate-TSAD has been a source of active concern. Wu and Keogh~\cite{wu2021current} document that several widely cited benchmarks exhibit triviality, mislabelling, and run-length bias that inflate reported scores; Lai et al.~\cite{lai2021revisiting} consolidate a formal taxonomy of point, contextual, and collective anomalies and propose evaluation protocols that disentangle these regimes. Together, these works argue for evaluating detectors on (i) labelled corpora authored by domain experts rather than by the model developers themselves and (ii) at least one dataset outside the training distribution. Hundman et al.~\cite{hundman2018detecting} provide the labelled NASA SMAP/MSL spacecraft telemetry corpus; Su et al.~\cite{su2019robust} provide the SMD server-machine corpus; Goh et al.~\cite{goh2017dataset} and Mathur and Tippenhauer~\cite{mathur2016swat} describe the SWaT water-treatment testbed; and Filonov, Lavrentyev, and Vorontsov~\cite{filonov2016multivariate} authored the predecessor of the SKAB testbed~\cite{skab_dataset} used in our cross-dataset experiment (\cref{sec:skab_cross_dataset}). We follow the cross-dataset protocol recommended by these critiques: the labelled evaluation is run on a corpus (SKAB) whose ground truth is exogenous to our model development pipeline and whose domain (industrial water circulation) is disjoint from the IAQ training distribution.

\begin{table}[H]
\centering
\caption{Comparison of the proposed framework with existing approaches across key dimensions.}
\label{table:positioning}
\begin{tabular}{lcccc}
\toprule
\textbf{Method} & \textbf{Uncertainty} & \textbf{Adaptive} & \textbf{Sensor} & \textbf{Feature} \\
 & \textbf{Quant.} & \textbf{Attention} & \textbf{Adaptation} & \textbf{Attribution} \\
\midrule
TCN-AE~\cite{thill_temporal_2021} & No & No & No & No \\
LSTM-ED~\cite{malhotra2016lstmed} & No & No & No & No \\
USAD~\cite{audibert2020usad} & No & No & No & No \\
Anomaly Transformer~\cite{xu2022anomalytransformer} & No & No & No & No \\
Stochastic TCN~\cite{ghimire_electricity_2025} & Yes & No & No & No \\
SDE-HNN~\cite{cui_sde-hnn_2025} & Yes & No & No & No \\
\textbf{This work} & \textbf{Yes} & \textbf{Yes} & \textbf{Yes} & \textbf{Yes} \\
\bottomrule
\end{tabular}
\end{table}